\section{Further Work}\label{sec:further-work}

The study of LLM personas is a burgeoning field with much work to be done, and we would love to see various improvements and extensions to this work. This section will not be exhaustive, but hopefully point towards low-hanging fruit that might help further validate and improve the assumptions and methodology.
\textbf{Improving the methodology for the supervised case}. Creating better LoRAs (i.e. ones with either more coverage of LLM behaviour, or with more independence allowing for cleaner combinations) would allow for more robust study of the relationship between trait space and personas. In addition, there would be benefit to further investigation of different ways to combine, and especially to invert, the trait LoRAs.
\textbf{Interpretability of the traits in weight space.} We would like to see exploration of where within a model different traits are shaped most, which could be done by further study into the weights themselves. Training LoRAs for the same trait across different models, or similar traits on the same model, could shed light into whether there is a consistent relationship between traits and weights, and whether continuity in one translates to continuity in the other.
\textbf{Further study into the relationship between trait space and personas.} Our work studying the traits we have has been limited to performing statistical analysis and visualising the flattened LoRA vectors. There could be a better (i.e. more predictive or more explainable) way of embedding this information, and we would be excited by work which seeks to discover this.
\textbf{More unsupervised exploration of personas.} From a safety perspective, understanding the non-human traits of LLMs is useful. Testing alternative methods to either naturally elicit diverse personas which exist within LLMs, or to cluster and label them, would allow better study of such behaviour. We would also be excited to see other, entirely different methods for the unsupervised exploration of personas.

